\documentclass[12pt,a4paper]{article}
\usepackage[romanian]{babel}
\usepackage[utf8]{inputenc}
\usepackage[T1]{fontenc}
\usepackage{geometry}
\usepackage{graphicx}
\usepackage{hyperref}
\usepackage{float}
\usepackage{amsmath}
\usepackage{listings}

\geometry{margin=2.5cm}

\lstset{
    basicstyle=\ttfamily\small,
    breaklines=true,
    columns=fullflexible
}

\title{Lucrarea de laborator nr. 4.2\\
Sistem de achizitie si conditionare a semnalului analogic}
\author{Nume Prenume\\Grupa}
\date{\today}

\begin{document}
\begin{titlepage}
\begin{center}
\vspace*{1cm}
{\Large Technical University of Moldova}\\[0.5cm]
{\Large Faculty of Computers, Informatics and Microelectronics}\\[2cm]
{\Huge\bfseries Laboratory Work No. 4.2}\\[1cm]
{\Large\bfseries Embedded Systems:\\
Analog Signal Acquisition and Conditioning\\
with Periodic STDIO Reporting}\\[3.5cm]
\begin{flushright}
\begin{tabular}{ll}
\textbf{Elaborated:} & \\
std. gr. FAF-231 & Burduja Adrian \\[1cm]
\textbf{Verified:} & \\
asist. univ. & Alexei Martîniuc \\ & \\
\end{tabular}
\end{flushright}
\vfill
{\Large Chișinău -- 2026}
\end{center}
\end{titlepage}
\newpage

%-------------------------------------------------
\section{Task Analysis and Technological Context}

\subsection{Domain Analysis and Technological Context}
Sensor-based embedded systems require more than raw data acquisition. In real applications,
sensor signals are subject to noise, measurement outliers, and physical range violations that
must be handled before the data can be used reliably. Signal conditioning pipelines — comprising
saturation, median filtering, and weighted averaging — are standard practice in industrial
monitoring, HVAC control, and robotics.

This work implements a periodic acquisition and conditioning pipeline for two sensors:
a DHT22 temperature and humidity sensor, and an HC-SR04 ultrasonic distance sensor.
The conditioned values are reported via a STDIO serial interface at a configurable interval.

\subsection{External Resources and Bibliographic References}
The following resources were used:
\begin{itemize}
    \item documentation of the Arduino Mega (ATmega2560) microcontroller;
    \item DHT22 sensor datasheet;
    \item HC-SR04 ultrasonic sensor datasheet;
    \item AVR-libc documentation (\texttt{dtostrf}, \texttt{fdev\_setup\_stream});
    \item educational materials provided within the course;
    \item Claude (Anthropic) for assistance with code generation and debugging.
\end{itemize}

\subsection{Analysis of Components and Used Resources}
The system uses:
\begin{itemize}
    \item Arduino Mega (ATmega2560) microcontroller;
    \item DHT22 sensor on pin A0 — temperature ($-40$\,°C to $80$\,°C) and humidity ($0$--$100$\,\%);
    \item HC-SR04 ultrasonic sensor on TRIG=7, ECHO=6 — distance ($2$--$400$\,cm);
    \item STDIO interface via serial port at 9600 baud;
    \item software resources: circular buffers, conditioning pipeline, millis-based scheduling.
\end{itemize}

%-------------------------------------------------
\section{Definition of Technical Requirements and Test Scenarios}

\subsection{Application Objectives}
The main objective is the design and implementation of an application that:
\begin{itemize}
    \item periodically acquires raw data from both sensors;
    \item applies a conditioning pipeline: saturation, median filter, weighted average;
    \item detects threshold crossings with hysteresis and persistence debouncing;
    \item reports all intermediate and final values via STDIO at a configurable interval.
\end{itemize}

\subsection{Functional and Non-Functional Requirements}

\textbf{Functional requirements:}
\begin{itemize}
    \item acquisition of temperature, humidity, and distance at a configurable interval;
    \item saturation of raw values to sensor physical limits;
    \item median filter over a window of 5 samples to suppress salt-and-pepper noise;
    \item weighted average giving higher weight to newer samples;
    \item threshold alerting with hysteresis ($\pm1$\,°C for temperature, $\pm2$\,cm for distance);
    \item persistence counter requiring 3 consecutive confirming samples before state change;
    \item structured STDIO report at 2000\,ms interval.
\end{itemize}

\textbf{Non-functional requirements:}
\begin{itemize}
    \item modular code separated into independent \texttt{.cpp}/\texttt{.h} files;
    \item no dynamic memory allocation;
    \item \texttt{printf}-based reporting via AVR-libc \texttt{fdev\_setup\_stream};
    \item float formatting via \texttt{dtostrf} due to AVR-libc \texttt{\%f} limitations.
\end{itemize}

\subsection{Test Scenarios and Validation Criteria}
\begin{itemize}
    \item temperature above $26$\,°C for 3 consecutive samples triggers ALERT state;
    \item temperature below $24$\,°C for 3 consecutive samples clears ALERT state;
    \item temperature in the dead zone [$24$\,°C, $26$\,°C] keeps current state unchanged;
    \item distance below $18$\,cm for 3 consecutive samples triggers TOO CLOSE alert;
    \item median output is stable under single-sample outliers;
    \item weighted output tracks trends faster than a simple moving average.
\end{itemize}

%-------------------------------------------------
\section{Architectural Design and Behavioral Modeling}

\subsection{System Components}
\begin{itemize}
    \item \textbf{Acquisition layer} — \texttt{dht\_read()}, \texttt{hcsr\_read()} read raw sensor values;
    \item \textbf{Conditioning layer} — \texttt{condition\_value()} applies saturation, median, weighted average;
    \item \textbf{Threshold layer} — \texttt{process\_threshold()}, \texttt{process\_distance()} implement hysteresis and persistence;
    \item \textbf{Reporting layer} — \texttt{report\_temperature()}, \texttt{report\_distance()} format and print via STDIO;
    \item \textbf{STDIO layer} — \texttt{stdio\_init()} binds \texttt{printf} to the hardware UART.
\end{itemize}

\subsection{HW--SW Layered Architecture}
The application is organized into the following layers:
\begin{itemize}
    \item hardware layer (DHT22, HC-SR04, UART);
    \item HAL layer (GPIO, \texttt{pulseIn}, \texttt{delayMicroseconds});
    \item conditioning layer (circular buffer, median sort, weighted sum);
    \item application layer (threshold logic, reporting).
\end{itemize}

\subsection{Conditioning Pipeline}
For each sensor value, the pipeline executes in order:

\begin{enumerate}
    \item \textbf{Saturation}: clamp raw value to physical sensor limits.
    \[
        v_{sat} = \min(\max(v_{raw},\, v_{min}),\, v_{max})
    \]

    \item \textbf{Median filter}: insert saturated value into a circular buffer of size $N=5$,
    sort a copy, return the middle element. Suppresses isolated outliers (salt-and-pepper noise).

    \item \textbf{Weighted average}: assign weight $w_i = i+1$ to the $i$-th oldest sample
    (newest sample gets weight $N$).
    \[
        \bar{v} = \frac{\sum_{i=0}^{N-1}(i+1)\cdot v_i}{\sum_{i=0}^{N-1}(i+1)}
    \]
\end{enumerate}

\subsection{Threshold and Hysteresis Logic}
To avoid alert oscillation near the threshold, two separate crossing levels are used:

\begin{itemize}
    \item Alert activates only when value exceeds $T_{high} = T_{threshold} + T_{hysteresis}$;
    \item Alert clears only when value drops below $T_{low} = T_{threshold} - T_{hysteresis}$;
    \item Values in the dead zone $[T_{low}, T_{high}]$ leave the state unchanged.
    \item State only changes after \texttt{PERSIST\_REQUIRED} consecutive confirming samples.
\end{itemize}

%-------------------------------------------------
\section{Electrical Schematic Design and Simulation}

\subsection{Connections and Component Sizing}
\begin{itemize}
    \item DHT22: data pin connected to A0, with a $10\,k\Omega$ pull-up resistor to $+5\,V$;
    \item HC-SR04: TRIG on pin 7 (output), ECHO on pin 6 (input), powered at $5\,V$;
    \item Serial output via USB at 9600 baud.
\end{itemize}

\subsection{Simulation}
The application was simulated in SimulIDE. Both sensors were configured with virtual
stimulus values to verify the conditioning pipeline and alert logic.

\subsection{Execution on Real Hardware}
Sample serial output:
\begin{lstlisting}
=== DHT22 Temperature Monitor (Mega) ===
------------------------------
Uptime (ms) : 2001
Temperature :  24.5 C
Humidity    :  58.0 %
Threshold   :  25.0 C  +/- 1.0 C
Persist     : 0/3
State       : NORMAL
Distance    :  35.2 cm
Dist Thresh :  20.0 cm
Dist State  : NORMAL
------------------------------
Uptime (ms) : 4002
Temperature :  26.3 C
Humidity    :  57.5 %
Threshold   :  25.0 C  +/- 1.0 C
Persist     : 1/3
State       : NORMAL
\end{lstlisting}

%-------------------------------------------------
\section{Modular Software Implementation}

\subsection{Project Structure}
\begin{lstlisting}
lab4-2.ino
main.cpp
conditioning.cpp
conditioning.h
dht22.cpp
dht22.h
hc-sr.cpp
hc-sr.h
stdio_layer.cpp
stdio_layer.h
\end{lstlisting}

\subsection{HW--SW Interface}
Sensor access is encapsulated behind \texttt{dht\_read()} and \texttt{hcsr\_read()},
which expose only the processed value to the rest of the application.
The conditioning function \texttt{condition\_value()} is fully generic and reused
for all three measured quantities.

\subsection{Scheduling}
No RTOS is used. Two \texttt{millis()}-based intervals control execution:
\begin{itemize}
    \item \texttt{ACQUIRE\_INTERVAL = 400\,ms} — sensor read and conditioning;
    \item \texttt{REPORT\_INTERVAL = 2000\,ms} — STDIO report.
\end{itemize}

\subsection{Comments and Documentation}
The code is structured into clearly named modules. Each module exposes only
the functions needed by other modules via its header file.

%-------------------------------------------------
\section{Simulation, Testing, and Validation}

\subsection{Testing}
Tests were performed by providing controlled sensor stimulus in SimulIDE and
verifying that the serial output reflected the expected conditioning and alert behavior.

\subsection{Results and Observations}
\begin{itemize}
    \item The median filter correctly rejected single-sample outliers without affecting steady-state output;
    \item The weighted average responded faster to temperature trends than a simple average;
    \item Hysteresis prevented alert flickering when the temperature hovered near the threshold;
    \item The persistence counter eliminated false alerts from transient spikes.
\end{itemize}

\subsection{Observed Issues and Fixes}
\begin{itemize}
    \item AVR-libc \texttt{printf} does not support \texttt{\%f} without special linker flags;
          \texttt{dtostrf()} was used instead;
    \item \texttt{LOW} and \texttt{HIGH} are macros in \texttt{Arduino.h}; local variables
          were renamed to \texttt{THRESH\_LOW} / \texttt{THRESH\_HIGH} to avoid silent
          macro substitution;
    \item Temperature and humidity conditioning windows were initially sharing the same
          position and count variables, causing corruption; they were separated into
          independent pairs (\texttt{g\_temp\_pos}/\texttt{g\_temp\_cnt} and
          \texttt{g\_humi\_pos}/\texttt{g\_humi\_cnt}).
\end{itemize}

%-------------------------------------------------
\section{Conclusions}

The developed application demonstrates a complete sensor acquisition and signal conditioning
pipeline on a resource-constrained AVR microcontroller. The modular architecture cleanly
separates acquisition, conditioning, threshold detection, and reporting. The use of a generic
\texttt{condition\_value()} function for all sensors reduces code duplication and simplifies
future extension to additional sensors. The solution is suitable as a foundation for more
complex embedded monitoring systems.

\newpage
\section{Annex}

\subsection{conditioning.h}
\begin{lstlisting}
#pragma once

#define MEDIAN_WINDOW 5

struct ConditionResults {
    float raw, saturated, median, weighted;
};

ConditionResults condition_value(
    float raw,
    float min_sat, float max_sat,
    float window[], int *window_pos, int *window_count
);
\end{lstlisting}

\subsection{conditioning.cpp}
\begin{lstlisting}
#include "conditioning.h"
#include <Arduino.h>

static float saturate(float val, float mn, float mx) {
    if (val < mn) return mn;
    if (val > mx) return mx;
    return val;
}

static float median(float window[], int count) {
    float buff[MEDIAN_WINDOW];
    memcpy(buff, window, count * sizeof(float));
    for (int i = 1; i < count; i++) {
        float key = buff[i];
        int j = i - 1;
        while (j >= 0 && buff[j] > key) { buff[j+1] = buff[j]; j--; }
        buff[j+1] = key;
    }
    return buff[count / 2];
}

static float weighted(float window[], int pos, int count) {
    float sum = 0, weight = 0;
    for (int i = 0; i < count; i++) {
        int   p = (pos + i) % count;
        float w = (float)(i + 1);
        sum    += window[p] * w;
        weight += w;
    }
    return sum / weight;
}

ConditionResults condition_value(
    float raw, float min_sat, float max_sat,
    float window[], int *window_pos, int *window_count
) {
    ConditionResults res = {};
    res.raw       = raw;
    res.saturated = saturate(raw, min_sat, max_sat);
    window[*window_pos] = res.saturated;
    *window_pos         = (*window_pos + 1) % MEDIAN_WINDOW;
    if (*window_count < MEDIAN_WINDOW) (*window_count)++;
    res.median   = median(window, *window_count);
    res.weighted = weighted(window, *window_pos, *window_count);
    return res;
}
\end{lstlisting}

\subsection{dht22.h}
\begin{lstlisting}
#pragma once
#include <DHT.h>
#include "conditioning.h"

#define DHT_PIN            A0
#define DHT_TYPE           DHT22

#define TEMP_THRESHOLD     25.0f
#define TEMP_HYSTERESIS     1.0f
#define PERSIST_REQUIRED    3

struct RawData {
    float temperature;
    float humidity;
    bool  valid;
};

extern DHT     dht;
extern RawData g_data;
extern bool    g_is_alert;
extern int     g_persist;

RawData          dht_read();
void             process_threshold();
void             report_temperature(unsigned long now);
ConditionResults dht_condition_temp();
ConditionResults dht_condition_humi();
\end{lstlisting}

\subsection{dht22.cpp}
\begin{lstlisting}
#include "dht22.h"
#include "conditioning.h"
#include <Arduino.h>

char buf[8];

DHT dht(DHT_PIN, DHT_TYPE);

RawData g_data     = { 0, 0, false };
bool    g_is_alert = false;
int     g_persist  = 0;

float g_temps[MEDIAN_WINDOW], g_humis[MEDIAN_WINDOW];
int   g_temp_pos = 0, g_temp_cnt = 0;
int   g_humi_pos = 0, g_humi_cnt = 0;

RawData dht_read() {
    g_data.temperature = dht.readTemperature();
    g_data.humidity    = dht.readHumidity();
    g_data.valid = !isnan(g_data.temperature) && !isnan(g_data.humidity);
    return g_data;
}

ConditionResults dht_condition_temp() {
    if (!g_data.valid) return ConditionResults{};
    return condition_value(g_data.temperature, -40, 80,
                           g_temps, &g_temp_pos, &g_temp_cnt);
}

ConditionResults dht_condition_humi() {
    if (!g_data.valid) return ConditionResults{};
    return condition_value(g_data.humidity, 0, 100,
                           g_humis, &g_humi_pos, &g_humi_cnt);
}

void process_threshold() {
    float temp = g_data.temperature;
    const float high_thresh = TEMP_THRESHOLD + TEMP_HYSTERESIS;
    const float low_thresh  = TEMP_THRESHOLD - TEMP_HYSTERESIS;
    if (!g_is_alert) {
        if (temp > high_thresh) {
            if (++g_persist >= PERSIST_REQUIRED)
                { g_is_alert = true;  g_persist = 0; }
        } else { g_persist = 0; }
    } else {
        if (temp < low_thresh) {
            if (++g_persist >= PERSIST_REQUIRED)
                { g_is_alert = false; g_persist = 0; }
        } else { g_persist = 0; }
    }
}

void report_temperature(unsigned long now) {
    printf("------------------------------\r\n");
    printf("Uptime (ms) : %lu\r\n", now);
    if (!g_data.valid) { printf("Sensor : READ ERROR\r\n"); return; }
    printf("Temperature : %s C\r\n", dtostrf(g_data.temperature, 5, 1, buf));
    printf("Humidity    : %s %%\r\n", dtostrf(g_data.humidity,   5, 1, buf));
    printf("Threshold   : %s C  +/-", dtostrf(TEMP_THRESHOLD,    5, 1, buf));
    printf("%s C\r\n",                 dtostrf(TEMP_HYSTERESIS,   4, 1, buf));
    printf("Persist     : %d/%d\r\n",  g_persist, PERSIST_REQUIRED);
    printf("State       : %s\r\n",
           g_is_alert ? "!!! ALERT !!!" : "NORMAL");
}
\end{lstlisting}

\subsection{hc-sr.h}
\begin{lstlisting}
#pragma once
#include "conditioning.h"

#define TRIG_PIN              7
#define ECHO_PIN              6

#define DIST_THRESHOLD       20.0f
#define DIST_HYSTERESIS       2.0f
#define DIST_PERSIST_REQUIRED 3

#define HCSR_DIST_MIN  2.0f
#define HCSR_DIST_MAX  400.0f

extern float g_distance;
extern bool  g_dist_alert;
extern int   g_dist_persist;

float            hcsr_read();
ConditionResults hcsr_condition_dist();
void             process_distance();
void             report_distance();
void             hcsr_init();
\end{lstlisting}

\subsection{hc-sr.cpp}
\begin{lstlisting}
#include "hc-sr.h"
#include "conditioning.h"
#include <Arduino.h>

char hcsr_buf[8];

float g_distance    = 0.0f;
bool  g_dist_alert  = false;
int   g_dist_persist = 0;

float g_dists[MEDIAN_WINDOW];
int   g_dist_pos = 0, g_dist_cnt = 0;

float hcsr_read() {
    digitalWrite(TRIG_PIN, LOW);
    delayMicroseconds(2);
    digitalWrite(TRIG_PIN, HIGH);
    delayMicroseconds(10);
    digitalWrite(TRIG_PIN, LOW);
    unsigned long duration = pulseIn(ECHO_PIN, HIGH, 30000UL);
    if (duration == 0) { g_distance = -1.0f; return g_distance; }
    g_distance = duration / 58.0f;
    return g_distance;
}

ConditionResults hcsr_condition_dist() {
    if (g_distance < 0) return ConditionResults{};
    return condition_value(g_distance, HCSR_DIST_MIN, HCSR_DIST_MAX,
                           g_dists, &g_dist_pos, &g_dist_cnt);
}

void process_distance() {
    float dist = g_distance;
    if (dist < 0) { g_dist_persist = 0; return; }
    const float THRESH_NEAR = DIST_THRESHOLD - DIST_HYSTERESIS;
    const float THRESH_FAR  = DIST_THRESHOLD + DIST_HYSTERESIS;
    if (dist < THRESH_NEAR && !g_dist_alert) {
        if (++g_dist_persist >= DIST_PERSIST_REQUIRED)
            { g_dist_alert = true;  g_dist_persist = 0; }
    } else if (dist > THRESH_FAR && g_dist_alert) {
        if (++g_dist_persist >= DIST_PERSIST_REQUIRED)
            { g_dist_alert = false; g_dist_persist = 0; }
    } else { g_dist_persist = 0; }
}

void hcsr_init() {
    pinMode(TRIG_PIN, OUTPUT);
    pinMode(ECHO_PIN, INPUT);
}

void report_distance() {
    printf("Distance    : %s cm\r\n",
           dtostrf(g_distance,     5, 1, hcsr_buf));
    printf("Dist Thresh : %s cm\r\n",
           dtostrf(DIST_THRESHOLD, 5, 1, hcsr_buf));
    printf("Dist State  : %s\r\n",
           g_dist_alert ? "!!! TOO CLOSE !!!" : "NORMAL");
}
\end{lstlisting}

\subsection{stdio\_layer.cpp}
\begin{lstlisting}
#include "stdio_layer.h"
#include <Arduino.h>

static int serial_putchar(char c, FILE *) {
    Serial.write(c);
    return 0;
}

static FILE g_stdout;

void stdio_init() {
    fdev_setup_stream(&g_stdout, serial_putchar, NULL, _FDEV_SETUP_WRITE);
    stdout = &g_stdout;
    Serial.begin(9600);
}
\end{lstlisting}

\subsection{main.cpp}
\begin{lstlisting}
#include <Arduino.h>
#include "conditioning.h"
#include "stdio_layer.h"
#include "dht22.h"
#include "hc-sr.h"

#define ACQUIRE_INTERVAL 400UL
#define REPORT_INTERVAL  2000UL

int main() {
    init();
    stdio_init();
    dht.begin();
    hcsr_init();
    printf("=== DHT22 Temperature Monitor (Mega) ===\r\n");

    unsigned long lastAcquire = 0;
    unsigned long lastReport  = 0;

    for (;;) {
        unsigned long now = millis();

        if (now - lastAcquire >= ACQUIRE_INTERVAL) {
            lastAcquire = now;
            RawData data = dht_read();
            float   dist = hcsr_read();
            if (data.valid) {
                g_data.temperature = dht_condition_temp().weighted;
                g_data.humidity    = dht_condition_humi().weighted;
                process_threshold();
            }
            if (dist >= 0) {
                g_distance = hcsr_condition_dist().weighted;
                process_distance();
            }
        }

        if (now - lastReport >= REPORT_INTERVAL) {
            lastReport = now;
            report_temperature(now);
            report_distance();
        }
    }

    return 0;
}
\end{lstlisting}

\end{document}